\documentclass[12pt,a4paper]{article}

\usepackage[a4paper,margin=1in]{geometry}
\usepackage[T1]{fontenc}
\usepackage[utf8]{inputenc}
\usepackage{lmodern}
\usepackage{xcolor}
\usepackage{booktabs}
\usepackage{longtable}
\usepackage{array}
\usepackage{enumitem}
\usepackage{tabularx}
\usepackage{titlesec}
\usepackage{fancyhdr}
\usepackage{hyperref}
\usepackage{lastpage}

\definecolor{iiitblue}{HTML}{123B63}
\definecolor{lightblue}{HTML}{EAF2F8}
\definecolor{lightgreen}{HTML}{EAF7EF}
\definecolor{lightorange}{HTML}{FFF4E5}
\definecolor{lightred}{HTML}{FDECEC}

\hypersetup{
    colorlinks=true,
    linkcolor=iiitblue,
    urlcolor=iiitblue,
    pdftitle={Pharmacy Inventory and Expiry Management System - Progress Report},
    pdfauthor={Shashi Prakash Yadav, Mohd Usham, Rishabh Kumar, Ritik Aryan}
}

\titleformat{\section}{\Large\bfseries\color{iiitblue}}{\thesection}{0.75em}{}
\titleformat{\subsection}{\large\bfseries\color{iiitblue}}{\thesubsection}{0.75em}{}
\pagestyle{fancy}
\fancyhf{}
\lhead{\small Pharmacy Inventory and Expiry Management System}
\rhead{\small Progress Report}
\cfoot{\small Page \thepage\ of \pageref{LastPage}}
\setlength{\headheight}{15pt}

\newcommand{\statusdone}{\textcolor{green!45!black}{\textbf{Completed}}}
\newcommand{\statusprogress}{\textcolor{orange!75!black}{\textbf{In progress}}}
\newcommand{\statusremaining}{\textcolor{red!70!black}{\textbf{Remaining}}}
\newcommand{\projecttitle}{Pharmacy Inventory \& Expiry Management System}

\begin{document}

\begin{titlepage}
\thispagestyle{empty}
\centering
\vspace*{1.0cm}
{\Large\bfseries INDIAN INSTITUTE OF\\[0.15cm] INFORMATION TECHNOLOGY, KALYANI\par}
\vspace{1.5cm}
{\color{iiitblue}\rule{\textwidth}{1.2pt}\par}
\vspace{0.6cm}
{\Huge\bfseries \projecttitle\par}
\vspace{0.35cm}
{\Large\bfseries Updated Project Progress Report\par}
\vspace{0.6cm}
{\color{iiitblue}\rule{\textwidth}{1.2pt}\par}
\vfill
\begin{tabular}{rl}
\textbf{Project Mentor:} & Dr. Sudeshna Mondal\\[0.25cm]
\textbf{Students:} & Shashi Prakash Yadav \; [CSE/23096/1144]\\
& Mohd Usham \; [ECE/23125/1173]\\
& Rishabh Kumar \; [ECE/23133/1181]\\
& Ritik Aryan \; [ECE/23134/1182]\\[0.25cm]
\textbf{Report date:} & 30 October 2026
\end{tabular}
\vfill
{\large\itshape Status of work completed to date and plan for completion\par}
\end{titlepage}

\tableofcontents
\newpage

\section*{Abstract}
\addcontentsline{toc}{section}{Abstract}
This report presents the current development status of the \projecttitle. The
system is being developed as a full-stack web application to reduce medicine
wastage, prevent stock-outs, improve batch-level traceability, and support
pharmacy staff in daily inventory operations. The implemented solution combines
a React and Vite frontend, a FastAPI and SQLAlchemy backend, a relational
database, role-based access control, expiry monitoring, FEFO dispensing,
inventory forecasting, procurement workflows, sales records, and exportable
reports.

The core application workflow is now implemented and connected end to end for
the main demonstration scenarios. The remaining work is primarily focused on
systematic testing, production deployment hardening, usability refinement,
documentation, and final evaluation using representative test data. This
document records both the completed work and the activities still required
before final submission.

\section{Project Objective and Scope}
\subsection{Problem Statement}
Manual registers and spreadsheets make it difficult to identify medicines that
are approaching expiry, select the correct batch for dispensing, maintain
accurate stock balances, and coordinate procurement. These limitations can lead
to avoidable wastage, stock-outs, delayed reporting, and weak audit trails.

\subsection{Objectives}
\begin{enumerate}[leftmargin=*]
    \item Maintain a structured catalogue of medicines, suppliers, batches, and
    pharmacy workspaces.
    \item Track stock quantities and expiry dates at batch level.
    \item Apply the First-Expiry-First-Out (FEFO) rule during dispensing.
    \item Generate low-stock, near-expiry, and demand-based replenishment
    recommendations.
    \item Record sales, purchases, returns, disposal, and operational audit
    events.
    \item Provide role-aware access for pharmacy administrators and pharmacists.
    \item Produce downloadable CSV and PDF reports for operational review.
\end{enumerate}

\section{Technology and Architecture}
\begin{table}[h]
\centering
\begin{tabularx}{\textwidth}{>{\bfseries}p{0.25\textwidth}X}
\toprule
Layer & Technologies and implementation \\
\midrule
Frontend & React 18, TypeScript, Vite, Chart.js, responsive CSS interface \\
Backend & Python, FastAPI, Pydantic schemas, SQLAlchemy ORM \\
Database & SQLite for local demonstration and PostgreSQL-ready configuration \\
Forecasting & XGBoost-based demand forecasting with a deterministic fallback \\
Reporting & CSV exports and PDF report endpoints \\
Communication & SMTP configuration for expiry and operational notifications \\
Deployment & Local development workflow with Docker Compose support for PostgreSQL \\
\bottomrule
\end{tabularx}
\caption{Current technology stack}
\end{table}

The application is organized as a multi-pharmacy platform. Records are linked
to the signed-in pharmacy workspace so that medicines, batches, procurement,
sales, and reports remain isolated between pharmacies. The backend creates and
updates the required database structures at startup and seeds synthetic
demonstration data for local evaluation.

\section{Work Completed to Date}
\subsection{Foundation and Access Control}
\begin{itemize}[leftmargin=*]
    \item Implemented a FastAPI backend with structured API routes, database
    sessions, ORM entities, validation schemas, and application configuration.
    \item Implemented pharmacy registration and multi-tenant workspace
    initialization.
    \item Implemented login and role-aware access for administrators and
    pharmacists.
    \item Added pharmacy-scoped access checks for operational records.
    \item Added a health endpoint and startup database compatibility handling.
\end{itemize}

\subsection{Medicine Catalogue and Inventory}
\begin{itemize}[leftmargin=*]
    \item Implemented medicine catalogue creation, editing, search, and
    category/reference data.
    \item Implemented batch-level stock records with batch number, quantity,
    expiry date, pricing, and disposal information.
    \item Added reorder levels, ideal stock values, low-stock indicators, and
    inventory summary metrics.
    \item Added CSV catalogue import with create-or-update behaviour by SKU.
    \item Added supplier records, purchase orders, and procurement requests.
\end{itemize}

\subsection{Expiry, FEFO, and Waste Accountability}
\begin{itemize}
    \item Implemented expiry risk zones for safe, approaching, critical, and
    expired stock.
    \item Implemented near-expiry monitoring and dashboard alerts.
    \item Implemented FEFO-assisted dispensing so the earliest suitable batch is
    selected first.
    \item Prevented expired batches from normal sale or dispensing workflows.
    \item Implemented supplier return and licensed biomedical-waste disposal
    records with disposal method, reference, date, and audit information.
    \item Added company recovery records for eligible expired stock, including
    recovery values, transferred quantities, and handover references.
\end{itemize}

\subsection{Sales, Transactions, and Reporting}
\begin{itemize}
    \item Implemented medicine dispensing with destination department and
    internal reference details.
    \item Implemented itemised sales invoices and sales ledger information.
    \item Implemented transaction and consumption history.
    \item Implemented dashboards for stock totals, alerts, daily transactions,
    sales, purchases, and forecast signals.
    \item Implemented downloadable CSV and PDF inventory reporting surfaces.
\end{itemize}

\subsection{Forecasting and Notifications}
\begin{itemize}
    \item Added 30-day, 60-day, and 90-day demand forecast views.
    \item Added XGBoost-based forecasting with a deterministic fallback when
    insufficient history or model conditions are available.
    \item Connected forecast and reorder recommendations to reviewable
    procurement requests.
    \item Added SMTP configuration and an email delivery check for expiry,
    stock, and supplier-return notifications.
\end{itemize}

\subsection{User Interface}
\begin{itemize}
    \item Implemented a multi-page React interface for login, dashboard,
    medicines, batches, expiry and waste, dispensing, sales, finance,
    imports, suppliers, transactions, departments, reports, and settings.
    \item Added charts and summary cards for operational visibility.
    \item Added responsive layouts, data tables, forms, alerts, and
    confirmation workflows.
    \item Added barcode/QR-style SKU or batch lookup support for dispensing
    workflows.
\end{itemize}

\section{Current Milestone Status}
\begin{longtable}{p{0.27\textwidth}p{0.17\textwidth}p{0.46\textwidth}}
\caption{Milestone status as of 30 October 2026}\\
\toprule
\textbf{Milestone} & \textbf{Status} & \textbf{Evidence in the implementation} \\
\midrule
\endfirsthead
\toprule
\textbf{Milestone} & \textbf{Status} & \textbf{Evidence in the implementation} \\
\midrule
\endhead
Project structure and backend foundation & \statusdone & FastAPI application, SQLAlchemy models, configuration, database session, and startup lifecycle. \\
Frontend application shell & \statusdone & React/Vite application with routed operational pages and reusable UI patterns. \\
Authentication and pharmacy workspaces & \statusdone & Login, administrator/pharmacist roles, pharmacy registration, and workspace scoping. \\
Medicine and batch inventory & \statusdone & CRUD workflows, batch expiry tracking, thresholds, pricing, and stock summaries. \\
Expiry and waste management & \statusdone & Risk zones, quarantine, supplier returns, recovery records, disposal references, and audit trail. \\
FEFO dispensing and sales & \statusdone & Earliest-expiry batch selection, dispensing records, invoices, and sales totals. \\
Procurement and suppliers & \statusdone & Suppliers, purchase orders, procurement requests, approval workflow, and status tracking. \\
Forecasting and alerts & \statusdone & Forecast views, reorder recommendations, low-stock and near-expiry alerts. \\
Exports and reporting & \statusdone & CSV/PDF reporting endpoints and reports page. \\
Automated testing and test evidence & \statusprogress & Functional scenarios still need to be executed systematically and documented. \\
Production deployment and security hardening & \statusremaining & Final environment configuration, secret review, deployment rehearsal, and access review remain. \\
Final documentation and presentation & \statusprogress & Technical documentation and final demonstration material remain to be consolidated. \\
\bottomrule
\end{longtable}

\section{Work Remaining}
\subsection{Testing and Validation}
\begin{enumerate}[leftmargin=*]
    \item Create a repeatable test dataset covering normal, low-stock,
    near-expiry, expired, returned, disposed, and recovered batches.
    \item Test all role permissions and confirm that one pharmacy cannot read or
    modify another pharmacy's records.
    \item Test FEFO selection with multiple batches and partial quantities.
    \item Test sales, purchases, returns, disposal, and recovery calculations
    against manually verified expected values.
    \item Test CSV import validation, duplicate SKUs, malformed rows, and
    empty-file behaviour.
    \item Test forecast behaviour with sufficient, insufficient, and unusual
    consumption history.
    \item Test PDF/CSV report contents and download behaviour.
    \item Record test cases, screenshots, defects, and fixes for the final
    evaluation chapter.
\end{enumerate}

\subsection{Deployment and Security Hardening}
\begin{enumerate}[leftmargin=*]
    \item Configure and validate PostgreSQL deployment using environment
    variables rather than local demonstration defaults.
    \item Review authentication, password handling, CORS, SMTP, and API error
    responses before demonstration or deployment.
    \item Confirm that real credentials, mailbox passwords, and private
    configuration values are not included in the repository.
    \item Add database backup and restore instructions and perform one
    recovery rehearsal.
    \item Run a final build and deployment rehearsal on a clean environment.
\end{enumerate}

\subsection{Usability and Documentation}
\begin{enumerate}[leftmargin=*]
    \item Conduct a short usability review with representative pharmacy
    workflows and refine labels, validation messages, and empty states.
    \item Add screenshots and verified workflow results to the final report.
    \item Complete API, installation, user-operation, and administrator
    documentation.
    \item Prepare the final presentation, demonstration script, and individual
    contribution summary.
\end{enumerate}

\section{Proposed Completion Plan}
\begin{table}[h]
\centering
\begin{tabularx}{\textwidth}{>{\bfseries}p{0.23\textwidth}p{0.18\textwidth}X}
\toprule
\textbf{Phase} & \textbf{Target} & \textbf{Activities} \\
\midrule
Validation & Week 1 & Execute functional, role, data-integrity, import, report, and forecast tests. \\
Hardening & Week 2 & Fix defects, review configuration and secrets, validate PostgreSQL deployment, and complete backup rehearsal. \\
Usability & Week 3 & Refine interface based on walkthroughs, capture screenshots, and finalize user instructions. \\
Submission & Week 4 & Complete final report, presentation, contribution summary, and final demonstration rehearsal. \\
\bottomrule
\end{tabularx}
\caption{Suggested remaining-work schedule}
\end{table}

\section{Conclusion}
The project has progressed from the initial interface and conceptual roadmap to
a connected pharmacy-management application with the main operational
workflows implemented. The strongest completed areas are batch-level
inventory, expiry accountability, FEFO dispensing, multi-pharmacy isolation,
procurement, sales, forecasting, and reporting.

The project is not yet considered final until the remaining validation,
deployment, security review, usability review, and documentation activities have
been completed. The next immediate priority is to run and document an
end-to-end test cycle using controlled synthetic data, followed by deployment
hardening and preparation of the final demonstration.

\section*{Declaration of Data Scope}
\addcontentsline{toc}{section}{Declaration of Data Scope}
The demonstration catalogue, suppliers, prices, stock quantities, and
consumption history are synthetic project data. They are intended for software
demonstration and testing only and do not represent real patient records,
licensed stock, or a guarantee of regulatory approval. Any operational
deployment must use verified local medicine, supplier, licensing, and disposal
information.

\vfill
\begin{center}
\rule{0.8\textwidth}{0.4pt}\\[0.2cm]
\small Prepared for mentor review\\
\small \projecttitle\ -- 30 October 2026
\end{center}

\end{document}
